\section{Interfaces and I/O Contracts}

This section defines the \textbf{structural} input/output contracts of the Engine
Core as a \textbf{diagnostic-only} kernel. The contracts are intentionally minimal:
they enable reproducible diagnostic evaluation while \textbf{explicitly forbidding}
control, optimization, adaptation, learning, or prescriptive semantics.

\subsection{Inputs (schemas, required fields)}

The Engine Core consumes exactly one evaluation request:
\[
Q := \langle S,\; \mathcal{I},\; \Theta,\; M \rangle
\]
where $S$ is the state snapshot, $\mathcal{I}$ the declared invariant set, $\Theta$
the declared threshold objects, and $M$ a minimal metadata envelope.

\paragraph{Required input components.}
\begin{itemize}
  \item \textbf{State snapshot $S$.} Immutable and read-only. Partiality is
        admissible \emph{only} if missing fields are explicitly represented in the
        observability descriptor $\mathcal{O}$.
  \item \textbf{Invariant set $\mathcal{I}$.} Explicit, versioned, and inspectable
        predicates $I_k : S \rightarrow \{0,1\}$.
  \item \textbf{Threshold objects $\Theta$.} Explicit, typed boundary parameters
        referenced by invariants. Thresholds are inputs only; they are never tuned,
        adapted, calibrated, or learned by the kernel.
  \item \textbf{Metadata envelope $M$.} Minimal contextual metadata such as
        \texttt{case\_id}, \texttt{engine\_version}, and stable identifiers of the
        input set. Metadata must not influence evaluation, classification, STOP,
        or provenance semantics.
\end{itemize}

\paragraph{Observability as explicit input.}
Unknown or indeterminate values must be represented explicitly (e.g.\ via
$\mathcal{O}$). The kernel must not replace missing observability by assumptions,
defaults, heuristics, estimation, or probabilistic completion.

\subsection{Outputs (artefacts, schemas)}

The Engine Core produces exactly one evaluation result:
\[
R := \langle \mathbf{v},\; \mathbf{c},\; \mathrm{STOP},\; P \rangle
\]
with:
\begin{itemize}
  \item $\mathbf{v}\in\{0,1\}^n$ — invariant satisfaction vector;
  \item $\mathbf{c}$ — finite structural classification label;
  \item $\mathrm{STOP}\in\{0,1\}$ — terminal STOP flag;
  \item $P$ — minimal \textbf{append-only} provenance record.
\end{itemize}

\paragraph{No prescriptive fields.}
Outputs must not contain recommendations, priorities, rankings, objectives,
targets, success criteria, or action proposals. The presence of any such field
invalidates OPS-1 conformance.

\paragraph{Provenance payload (minimal).}
At minimum, the provenance record binds outputs to inputs via stable identifiers:
\[
P := \langle h(S),\, h(\mathcal{I}),\, h(\Theta),\, \mathbf{v},\, \mathbf{c},\, \mathrm{STOP} \rangle
\]
External systems may attach additional identifiers (e.g.\ run IDs or build hashes)
\emph{only} if they introduce no new semantics and do not alter classification,
STOP meaning, admissibility boundaries, or provenance identity.

\subsection{Error surfaces (typed failure outputs)}

The Engine Core strictly distinguishes \textbf{terminal STOP} from
\textbf{tooling or implementation errors}.

\paragraph{STOP (kernel-level).}
STOP is a first-class, terminal boundary of diagnostic admissibility
(see Section~\ref{sec:stop}). It is neither a recoverable error, nor a warning,
nor advice.

\paragraph{Implementation errors (non-kernel).}
Failures such as parsing errors, missing files, corrupted input containers, or
environment misconfiguration are \emph{external} to Engine Core semantics. They may
be represented by surrounding tooling as typed errors, but:
\begin{itemize}
  \item they must not be conflated with STOP;
  \item they must not introduce recovery, retry, or fallback logic into the kernel;
  \item they must not affect invariant evaluation, classification, or STOP semantics.
\end{itemize}

The kernel defines $R$ only for admissible inputs. Input admission, error handling,
and tooling reliability are external concerns.

\subsection{Extension points and compatibility boundary (OPS-1)}

\paragraph{Allowed extension points (closed list).}
The following extension points are \textbf{explicitly allowed}. This list is
\textbf{closed}: any extension not listed here is non-admissible under OPS-1.
\begin{itemize}
  \item \textbf{Presentation or rendering layers} (cosmetic-only; no semantic effect).
  \item \textbf{Non-normative reporting views} (read-only; no decision effect).
  \item \textbf{Evidence import adapters} (strict: no semantic change; full provenance).
  \item \textbf{Additional verification gates} (stricter allowed; never weaker).
  \item \textbf{Witness packaging} (bundling artefacts; no new meaning).
\end{itemize}

\paragraph{Semantic neutrality constraint.}
Every allowed extension is admissible \emph{only} if it introduces
\textbf{no new semantic degrees of freedom}. It must not alter classification,
STOP semantics, provenance identity, or admissibility boundaries.

\paragraph{Binary conformance.}
Conformance is \textbf{binary}: conformant or \emph{non-ESTRA}. Any layer that
introduces a forbidden class (see Limits and Non-Goals) is non-admissible,
regardless of tooling quality or intent.

\paragraph{Toolkit / CLI / GUI boundary (OPS-2 as consumer only).}
Any toolkit surface (CLI/GUI/scripts/integrations) is outside OPS-1 and belongs to
OPS-2. OPS-2 is an \textbf{exchangeable consumer} of the OPS-1 contract: it consumes
$Q$ and renders or transports $R$, but it must not define, reinterpret, extend, or
operationalize OPS-1 semantics.

OPS-2 variation is admissible \emph{only} within the closed extension space above.
No interface layer may weaken invariants or gates, nor introduce control semantics,
feedback loops, learning, heuristic shortcuts, or STOP bypass.

\paragraph{Auto-disqualification by semantic contamination.}
If a toolkit layer (or its documentation or positioning) implies that the Engine
Core is a ``tool'', ``method'', ``framework'', ``best practice'', ``playbook'', or
``how-to'', that representation is non-conformant by definition.
If a toolkit layer induces or implies ``you should do X'' from $R$ or from STOP,
it introduces prescriptive semantics and is therefore \textbf{non-ESTRA}.

\subsection{Versioning and compatibility contract}

\paragraph{Version identifiers.}
The Engine Core must expose a version identifier (e.g.\ content hash or semantic
version) via metadata $M$ and/or provenance $P$.

\paragraph{Compatibility rule (minimal).}
Results are comparable as kernel outputs \emph{only if}:
\begin{itemize}
  \item $h(\mathcal{I})$ is identical;
  \item $h(\Theta)$ is identical;
  \item the Engine Core version identifier is identical.
\end{itemize}

If any of these differ, comparison may occur externally, but it is cross-version
analysis and not a kernel property.

\paragraph{No implicit upgrades.}
The kernel must not silently reinterpret invariants, thresholds, classifications,
or STOP semantics across versions. Any change requires an explicit version change
and updated identifiers.
